\section{n4. Exact crossing-surface conductor criterion; mixed-characteristic failures and a normal-component surface classification.}
\textit{Supports: G331, n1, n2, n3.}

There is an EXACT obstacle-by-obstacle existence criterion on the normal-crossing surfaces of n2, and a sharp universal CE criterion for excellent reduced surfaces with normal irreducible components. These are new scoped results, using ENTIRE qualified G331 and own n1--n3.

First let R=k[x,y,z]\_(x,y,z)/(xy) over any field, or R=V[[x,y]]/(xy) for a complete discrete valuation ring V, with primary gauge m=(x,y,z) or m=(x,y,pi). In the mixed-characteristic case the finite normalization is V[[x]] times V[[y]], R is their fibre product over V, and conductor is(x,y). These follow by splitting formal series into their common constant-in-x,y part and the two separated remainders. Both factors are regular excellent normal surfaces, with equal fraction fields componentwise; the conductor test by idempotents is unchanged. The two conductor-boundary points map to one o. G331 supplies the full quotient and positive C/D descent.

For ANY finite continuous signed obstacle f on X, write f\_i for its normalization restrictions and H\_i=P\_i(f\_i) for their greatest continuous geometric envelopes, which exist by n1's arbitrary-characteristic normal-surface theorem. Let o\_i be the conductor point upstairs and a=min(H\_1(o\_1),H\_2(o\_2)). Then the ENTIRE raw downstairs C/D supremum is
p(rho(w))=H\_i(w) when w is on component i and w(x\_i)\textless{}infinity; p(o)=a,
where x\_1=x,x\_2=y. It has a greatest C/D minorant precisely when H\_1(o\_1)=H\_2(o\_2). In that case p is continuous and belongs to geometric C. Otherwise p is not USC and neither greatest minorant exists. No fibrewise minimum is silently claimed to be psh.

Proof of the lower bound, including actual positivity. G331's normal zero dichotomy implies that H\_i(o\_i)=0 forces H\_i identically0. If a=0, both envelopes are0, f\_i\textgreater{}=0 and p=0. If a\textless{}0, for H\_i(o\_i)\textless{}0 put B\_i=lambda\_i H\_i with lambda\_i=a/H\_i(o\_i)\textgreater{}=1; for H\_i identically0 put B\_i equal to constant a. In either case B\_i is in geometric C, B\_i\textless{}=H\_i\textless{}=f\_i, and B\_i(o\_i)=a. For n\textgreater{}=1 put
G\_(i,n)=max(B\_i,H\_i-w(x\_i)/n),
with the second term bottom where x\_i is killed. This is a geometric C function: choose an integer N\textgreater{}=n||H\_i-B\_i||. Since the gauge is1, w((x\_i)+m\_i\textasciicircum{}N)=min(w(x\_i),N), and
G\_(i,n)=max(B\_i,H\_i-(1/n)w((x\_i)+m\_i\textasciicircum{}N)).
The right side belongs to C by the granted addition, scaling and maximum operations and the ACTUAL negative effective ample primary divisor of (x\_i)+m\_i\textasciicircum{}N. Equality holds even at infinite value because H\_i-N/n\textless{}=B\_i everywhere. Thus G\_(i,n) is finite continuous, equals a at o\_i, and is below f\_i. The compatible tuple descends WITH GEOMETRIC POSITIVITY by G331. It increases to H\_i at every finite-x\_i point and keeps value a at the conductor. This proves the claimed raw lower bound. Any target D competitor pulls to a source D minorant, hence is\textless{}=H\_i, and its shared value is\textless{}=both H\_i(o\_i). This proves the upper bound on ALL supports.

When the two conductor values agree, the compatible tuple H\_i descends to a continuous geometric envelope. If, say, H\_1(o\_1)\textgreater{}a, generic monomial w\_r(x)=r,w\_r(z)=1 tends on every element to o\_1; in the DVR case use coefficient-order minima w\_r(sum a\_j x\textasciicircum{}j)=min\_j(ord\_pi(a\_j)+rj), w\_r(pi)=1. The associated graded domain proves multiplicativity, all units have value0, and as r tends to infinity the value tends to ord\_pi(a\_0), or infinity if a\_0=0. Continuity of H\_1 makes p(w\_r) tend to H\_1(o\_1)\textgreater{}p(o). Thus p is not USC and cannot be a greatest D or C minorant. For finite q\textgreater{}=2 the same proof applies to an EXCELLENT local fibre product of excellent regular local SURFACE rings with surjective maps to the SAME marked DVR, nonzero principal kernels (x\_i), and the product as its finite normalization. Its gauge must extend to primary gauges on every factor; each DVR boundary is normalized by that same quotient gauge. Then a is the minimum of all q conductor values, use the same B\_i and G\_(i,n) with (x\_i)+a\_i\textasciicircum{}N, and require equality of ALL conductor values. At an incompatible value use G331 divisorial density to obtain a convergent NET in the finite-kernel-value locus; no sequence assumption is necessary. These exact assumptions hold for k[x\_1,...,x\_q,z]\_0/(x\_i x\_j:i!=j) and V[[x\_1,...,x\_q]]/(x\_i x\_j:i!=j), whose fibre products and conductors follow by separated-term expansion. No higher-dimensional conductor min-field is claimed geometric by this argument.

The explicit f,F,h\_n and canonical filtration in n2 also work on V[[x,y]]/(xy), replacing z by pi. The branch2 generic m-order is min\_(j)(ord\_pi(a\_j)+j) for a formal y-series. A generator with conductor order e has a nonzero pi\textasciicircum{}e coefficient in its y=0 residue, so its generic order is\textless{}=e. Every arbitrary-support normalized semivaluation dominates ordinary m-order because m\textasciicircum{}e is finitely generated and units have value0. Thus the same geometric inequalities, actual primary positivity and exact canonical ideals hold; the every-degree error is1/2 and there is no greatest C/D minorant. This supplies a nonnormal mixed-characteristic failure while n1's mixed cusp remains a CE example.

For EVERY d\textgreater{}=2 and field k, the n2 obstruction and exact filtration extend to R\_d=k[x,y,z\_1,...,z\_(d-1)]\_0/(xy), m its maximal ideal. The conductor quotient is the z-ring; all its supports are retained. Set f=-1/2 on branch1 and -1/2-1/(2w(y)) on branch2, and F=(1/2)v((y)+m\textasciicircum{}2)-(3/2)v(m). Both agree on the ENTIRE infinite-conductor stratum. The branch2 ordinary m-order valuation and the conductor point killing x,y,z\_j(j\textgreater{}=2) with z\_1=1 give v1(I)\textless{}=v\_o(I). The h\_n from n2 have the identical all-support branch formulas. Every branch2 point is bounded by-1; every finite-x point on branch1 reaches-1/2; every conductor point is bounded by-1 because the branch2 m-order is dominated there. Hence raw p is-1/2 on the finite-x first-component locus and-1 otherwise. The canonical ideals are
A\_m=(z\_1,...,z\_(d-1))\textasciicircum{}m+x(x,z\_1,...,z\_(d-1))\textasciicircum{}(ceil(m/2)-1)+y(y,z\_1,...,z\_(d-1))\textasciicircum{}(m-1).
Their profiles have the same expression with t=min\_j v(z\_j), the same pointwise all-degree limit and the same supremum error1/2. Monomial rays x=r, all z\_j=1 converge to the conductor semivaluation killing x; units and all polynomials obey pointwise convergence. These failures use no higher-dimensional source CE theorem. An analogous extension to more branches uses separated x\_i\textasciicircum{}2 m\textasciicircum{}n on all branches except branch2; no new universal positive theorem follows.

Universal criterion with normal components. Let R be an excellent reduced local ring of dimension at most TWO, whose finitely many irreducible component rings R/p\_i are normal local domains. Then ALL finite continuous signed obstacles have continuous geometric C/D envelopes iff every pairwise intersection has dimension ZERO: dim R/(p\_i+p\_j)=0 for all i!=j. Necessity is n3; it even supplies signed primary-model failures. For sufficiency, S=product R/p\_i is the finite normalization. A semivaluation support strictly below m cannot contain two minimal primes, since their sum has radical m. Its unique component quotient fixes its entire lift; within a normal component normalization is the identity. Hence the FULL boundary map is bijective, including any conductor support. By G331 and n1 it is an order/isometric C/D equivalence, and the normal component CE theorems, direct DVR and empty cases give CE downstairs. Equivalently in this setting S/R is supported at the closed point: on the punctured base only one normal component is present. The conductor therefore contains a power of m; it is the unit ideal in the normal case and a proper primary ideal in the genuinely nonnormal case. This is a complete classification in the stated surface sector, not for irreducible nonnormal germs or arbitrary singular higher dimension.

More generally, at ANY finite normalization for which source CE holds, constancy of the source envelope on EVERY target fibre is sufficient for continuous positive descent and equality of envelopes, by G331's full positive descent and the source upper bound on every pullback competitor. Necessity of that condition for arbitrary targets is not asserted. For the explicit surface DVR gluing it is necessary, by the proved formula above. Every gauge transport is homogeneous and uses bounded positive denominators; the pointwise formula and existence condition transform together. Reduced-curve min(f,0), the literal radial distinction, every-support conventions, nilpotent reduction and empty boundaries remain unchanged. Artifacts n4.md,n4.tex; rerun python diagnostics.py; python deliver.py; python check\_delivery.py. These are complete scoped mathematical deductions, not entire ORIGINAL P12.3 closure.

\subsection*{Exact surface gluing and classification}
\[
 a=\min_i H_i(o_i),\qquad
 p(\rho w)=H_i(w)\ (w(x_i)<\infty),\qquad p(o)=a.
\]
For $B_i(o_i)=a$ and $B_i\le H_i$ the positive gluing functions are
\[
 G_{i,n}=\max\{B_i,H_i-w(x_i)/n\}
 =\max\{B_i,H_i-v((x_i)+\mathfrak a_i^N)/n\}
\]
whenever $N\ge n\|H_i-B_i\|_\infty$.
\[
 \text{greatest C/D minorant exists}\ \Longleftrightarrow\ H_i(o_i)=a\ \forall i.
\]
For excellent reduced surfaces with normal irreducible components,
\[
 \operatorname{CE}(R)\ \Longleftrightarrow\
 \dim R/(\mathfrak p_i+\mathfrak p_j)=0\quad(i\ne j).
\]

